% !TEX root = ../../tir_monograph_v12.tex
\chapter[Half Coordinate, Binary Information and Phase Circle]{The Half Coordinate, Binary Information, and the Phase Circle}
\label{ch:v12_half_ln2}

\section{Binary probability coordinate}

The primitive relational pair from Chapter~\ref{ch:v12_first_distinction} admits the normalized coordinate
\begin{equation}
\mathbf p(u)=(1-u,u),
\qquad
0\le u\le1.
\end{equation}
Pole exchange is
\begin{equation}
J(u)=1-u.
\end{equation}
Its unique fixed point is
\begin{equation}
\boxed{u_\star=\frac12.}
\label{eq:v12_half_point}
\end{equation}
Hence
\begin{equation}
p_N=p_S=\frac12,
\qquad
\boxed{H_2(1/2)=\ln2.}
\label{eq:v12_half_entropy}
\end{equation}
This scalar branch is fixed before flavour, particle, or spatial identifications enter.
\vTwelveStatus{B}{--}{PASS}

\section{Coherent phase fibre over the half seam}

Restoring the relative phase gives the balanced two-state family
\begin{equation}
\boxed{
|\psi_{1/2}(\varphi)\rangle
=
\frac{|N\rangle+e^{i\varphi}|S\rangle}{\sqrt2},
\qquad
\varphi\in\mathbb R/2\pi\mathbb Z.
}
\label{eq:v12_half_fiber_state}
\end{equation}
The scalar probability projection forgets $\varphi$, so
\begin{equation}
\boxed{
\pi^{-1}\!\left(\frac12,\frac12\right)
\cong
U(1)
\cong
S^1.
}
\label{eq:v12_half_fiber}
\end{equation}

The foundational $U(1)\cong S^1$ carrier is owned by \texttt{TIR\_RELATIONAL\_PHASE\_LAGRANGIAN\_CORE\_V0\_1.md}. The balanced family above is its exact half-seam realization on the two-state carrier.

\section{Pole exchange on the phase circle}

The exchange operator
\begin{equation}
X|N\rangle=|S\rangle,
\qquad
X|S\rangle=|N\rangle
\end{equation}
acts projectively as
\begin{equation}
\boxed{\varphi\mapsto-\varphi\pmod{2\pi}.}
\label{eq:v12_half_phase_exchange}
\end{equation}
Its fixed phases are
\begin{equation}
\boxed{\varphi\in\{0,\pi\}.}
\label{eq:v12_half_fixed_phases}
\end{equation}

\section{From the phase circle and the polar roles to the sphere}

The relational dyad already supplies two distinguished endpoint roles, denoted $N$ and $S$. The phase fibre supplies the equatorial circle $S^1$. The canonical theorem source \texttt{TIR\_LAGRANGIAN\_BLOCH\_SELECTION\_V0\_1.md} identifies the endpoint quotient with the standard topological suspension:
\begin{equation}
\boxed{\Sigma S^1\cong S^2.}
\label{eq:v12_suspension_s1_s2}
\end{equation}
Explicitly,
\begin{equation}
S^2
=
\left(S^1\times[-1,1]\right)\Big/
\left(
S^1\times\{-1\}\sim S,\;
S^1\times\{+1\}\sim N
\right).
\end{equation}

Therefore the sphere is obtained without assuming a prior $SO(3)$ action:
\begin{equation}
\boxed{
S^1_{\rm phase}+\{N,S\}_{R}
\longrightarrow
\Sigma S^1
\cong
S^2.
}
\label{eq:v12_phase_poles_sphere}
\end{equation}

\section{Euler-characteristic closure}

For a finite CW complex $X$, suspension obeys
\begin{equation}
\chi(\Sigma X)=2-\chi(X).
\end{equation}
Since
\begin{equation}
\chi(S^1)=0,
\end{equation}
one obtains
\begin{equation}
\boxed{\chi(\Sigma S^1)=2=\chi(S^2).}
\label{eq:v12_euler_characteristic_sphere}
\end{equation}

This Euler-characteristic relation is topological. It is distinct from Euler's complex phase identity $e^{i\pi}=-1$, which enters only after the Berry connection is available in Chapter~\ref{ch:v12_complex_carrier}.

\section{Explicit spherical realization}

Introduce latitude through
\begin{equation}
u=\sin^2\frac{\vartheta}{2},
\qquad
1-u=\cos^2\frac{\vartheta}{2}.
\end{equation}
Then
\begin{equation}
\boxed{
\mathbf r(u,\varphi)
=
\left(
2\sqrt{u(1-u)}\cos\varphi,\,
2\sqrt{u(1-u)}\sin\varphi,\,
1-2u
\right)
}
\label{eq:v12_half_bloch_map}
\end{equation}
satisfies identically
\begin{equation}
|\mathbf r|^2
=
4u(1-u)+(1-2u)^2
=
1.
\end{equation}
Thus it realizes the suspension sphere explicitly.

At the balanced point,
\begin{equation}
\boxed{
\mathbf r\!\left(\frac12,\varphi\right)
=
(\cos\varphi,\sin\varphi,0),
}
\label{eq:v12_half_equator}
\end{equation}
so the half-seam fibre is the equator, while $u=0$ and $u=1$ collapse the phase circle to the two poles.

Hence
\begin{equation}
\boxed{
N
\longleftrightarrow
S^1_{\rm half}
\longleftrightarrow
S
}
\end{equation}
is not a pictorial analogy but the explicit coordinate realization of $\Sigma S^1$.

\section{Phase closure and arithmetic indices}

For a closed phase loop,
\begin{equation}
\boxed{
n=\frac1{2\pi}\oint d\varphi\in\mathbb Z.
}
\label{eq:v12_winding}
\end{equation}
Equivalently, with $z=e^{i\varphi}$,
\begin{equation}
z^n=1.
\end{equation}
Thus
\begin{equation}
\boxed{
S^1
\longleftrightarrow
U(1)
\longrightarrow
\text{winding/degree}
\longrightarrow
n\in\mathbb Z.
}
\end{equation}
Positive closure magnitude gives the corresponding natural-number index.

\section{Chapter output}

The chapter exports
\begin{equation}
\boxed{
\{N,S\}
\to
\frac12
\to
\ln2,
\qquad
\frac12
\to
U(1)\cong S^1,
\qquad
S^1+\{N,S\}
\to
\Sigma S^1\cong S^2.
}
\label{eq:v12_ch2_output}
\end{equation}

Chapter~\ref{ch:v12_complex_carrier} now supplies the complex/projective realization $S^2\cong\mathbb{CP}^1$, the Fubini--Study metric and Berry curvature. Only after that geometric structure exists is the Euler--Berry spin-selection theorem invoked.

\section{Ownership boundary}

Chapter 2 owns the half coordinate, $\ln2$, the coherent $U(1)\cong S^1$ fibre, and the topological $S^2$ closure. Chapter 3 owns the complex/projective realization, Fubini--Study/Berry geometry, and the spinorial closure. Chapter 9 remains the sole v12 publication owner of
\[
\kappa=\frac{\ln2}{24\pi}.
\]
